\begin{proofbox}
\textbf{方法一：平方差配合和差，只平方一次（优先）。}
设焦点 $F_1(-c,0),F_2(c,0)$，定义中的距离差为 $2a$，且 $c>a>0$。
右支点 $P(x,y)$ 到两焦点距离记为 $r_1,r_2$。
\begin{center}\begin{tikzpicture}[font=\small,>=stealth]
\draw[black!30,rounded corners=2pt] (0,0) rectangle (14.4,8.2);
\node[anchor=west,font=\small\bfseries] at (.5,7.72) {右支的两条焦半径：先用平方差求和，再由和差求距离};

\begin{scope}[shift={(6.2,3.1)},scale=.58]
\draw[black!35] (-6,0)--(7,0);
\draw[blue,very thick,domain=-2.5:5.5,samples=121,variable=\t] plot ({3*sqrt(1+\t*\t/16)},\t);
\draw[blue,very thick,domain=-2.5:5.5,samples=121,variable=\t] plot ({-3*sqrt(1+\t*\t/16)},\t);
\draw[orange!85!black,very thick] (-5,0)--(5,5.3333333333);
\draw[violet,very thick] (5,0)--(5,5.3333333333);
\fill (-5,0) circle (2pt) node[below=5pt] {$F_1(-5,0)$};
\fill (5,0) circle (2pt) node[below=5pt] {$F_2(5,0)$};
\fill (5,5.3333333333) circle (2pt) node[above left=5pt] {$P(5,16/3)$};
\node[above left=4pt] at (0,2.6666666667) {$r_1=34/3$};
\node[right=7pt] at (5,2.6) {$r_2=16/3$};
\end{scope}

\node[anchor=west] at (.7,.55) {$r_1-r_2=6=2a$，$r_1^2-r_2^2=4cx$；故 $r_1+r_2=2cx/a$};
\end{tikzpicture}\end{center}

右支 $x>0$，故 $r_1>r_2$，由第一定义及距离平方差，
\[
 r_1-r_2=2a,\quad r_1^2-r_2^2=4cx
 \quad\Longrightarrow\quad r_1+r_2=\frac{2cx}{a}.
\]
和差求出 $r_2=cx/a-a$，再用距离公式平方一次：
\[
 (x-c)^2+y^2=(cx/a-a)^2.
\]
展开消去一次项，得到
\[
 (c^2-a^2)x^2-a^2y^2=a^2(c^2-a^2),\qquad
 \frac{x^2}{a^2}-\frac{y^2}{c^2-a^2}=1.
\]
与标准方程比较得 $b^2=c^2-a^2$，即 $c^2=a^2+b^2$。
左支由关于 $y$ 轴的对称性得到同一方程；纵向情形交换坐标即可。
这里是从定义推出关系，下面再从方程验证候选焦点，补足反向检查。
\end{proofbox}
\clearpage
\begin{proofbox}
\textbf{方法二：从方程构造焦点，分支取正根。}
给定 $x^2/a^2-y^2/b^2=1$，先构造 $c_0=\sqrt{a^2+b^2}>a$，候选焦点 $(\pm c_0,0)$。
方程给出 $|x|\ge a$ 及 $y^2=b^2(x^2/a^2-1)$，于是直接展开可验
\[
 (x\pm c_0)^2+y^2=(c_0x/a\pm a)^2.
\]
右支 $x\ge a$ 时，$c_0x/a-a\ge c_0-a>0$，两距离为
$c_0x/a+a$ 和 $c_0x/a-a$，差为 $2a$。
左支 $x\le-a$ 时，两括号都为负，取绝对值后两距离为
$-c_0x/a-a$ 和 $-c_0x/a+a$，差的绝对值仍为 $2a$。
因此这两个定点确可作为焦点，半焦距 $c=c_0$，得到 $c^2=a^2+b^2$。

\medskip
\textbf{开口角的直观延伸。} 渐近线由二次项相消给出 $y=\pm(b/a)x$；
两条指向右方的射线夹角为 $2\arctan(b/a)$。
\begin{center}\begin{tikzpicture}[font=\small,>=stealth]
\draw[black!30,rounded corners=2pt] (0,0) rectangle (14.4,8.2);
\node[anchor=west,font=\small\bfseries] at (.5,7.72) {固定实半轴 $a=3$：$b/a$ 越大，两条渐近线张角越大};
\begin{scope}[shift={(2.4,4)},scale=.24]
\draw[black!35] (-7,0)--(7,0);
\draw[black!45,dashed] (-5,-1.6666666666666667)--(5,1.6666666666666667);
\draw[black!45,dashed] (-5,1.6666666666666667)--(5,-1.6666666666666667);
\draw[blue,very thick,domain=-2:2,samples=101,variable=\t] plot ({3*sqrt(1+\t*\t)},{1*\t});
\draw[blue,very thick,domain=-2:2,samples=101,variable=\t] plot ({-3*sqrt(1+\t*\t)},{1*\t});
\fill (3.1622776602,0) circle (4pt);\fill (-3.1622776602,0) circle (4pt);
\end{scope}
\node at (2.4,1.2) {$b=1,\quad b/a=1/3$};
\begin{scope}[shift={(7.2,4)},scale=.24]
\draw[black!35] (-7,0)--(7,0);
\draw[black!45,dashed] (-5,-5.0)--(5,5.0);
\draw[black!45,dashed] (-5,5.0)--(5,-5.0);
\draw[blue,very thick,domain=-2:2,samples=101,variable=\t] plot ({3*sqrt(1+\t*\t)},{3*\t});
\draw[blue,very thick,domain=-2:2,samples=101,variable=\t] plot ({-3*sqrt(1+\t*\t)},{3*\t});
\fill (4.2426406871,0) circle (4pt);\fill (-4.2426406871,0) circle (4pt);
\end{scope}
\node at (7.2,1.2) {$b=3,\quad b/a=3/3$};
\begin{scope}[shift={(12,4)},scale=.24]
\draw[black!35] (-7,0)--(7,0);
\draw[black!45,dashed] (-5,-8.333333333333334)--(5,8.333333333333334);
\draw[black!45,dashed] (-5,8.333333333333334)--(5,-8.333333333333334);
\draw[blue,very thick,domain=-2:2,samples=101,variable=\t] plot ({3*sqrt(1+\t*\t)},{5*\t});
\draw[blue,very thick,domain=-2:2,samples=101,variable=\t] plot ({-3*sqrt(1+\t*\t)},{5*\t});
\fill (5.8309518948,0) circle (4pt);\fill (-5.8309518948,0) circle (4pt);
\end{scope}
\node at (12,1.2) {$b=5,\quad b/a=5/3$};

\node[anchor=west] at (.7,.55) {同尺度；虚线 $y=\pm(b/a)x$；黑点为焦点，$e=\sqrt{1+(b/a)^2}$};
\end{tikzpicture}\end{center}

所以这里的“开口越开阔”具体指\textbf{渐近线夹角越大}，由比值 $b/a$ 决定；
不能仅凭某一个绝对长度比较不同缩放的图。
\end{proofbox}
\clearpage
